\documentclass[12pt]{article}
\usepackage[a4paper,margin=2.2cm]{geometry}
\usepackage{amsmath,amssymb,amsthm,mathtools}
\usepackage{hyperref}
\usepackage{cleveref}

% 中文（xelatex）
\usepackage{fontspec}
\usepackage{xeCJK}

\usepackage{fontspec}
\usepackage{xeCJK}
\setCJKmainfont{Noto Serif CJK SC}

\title{Numerical QR}
\date{}
\begin{document}
\maketitle
\begin{flushleft}
    1: 几乎线性相关的矩阵 \\
    取: \\
    \begin{align}
        A_{\varepsilon} = \begin{bmatrix}
                              1           & 1           & 1           \\
                              \varepsilon & 0           & 0           \\
                              0           & \varepsilon & 0           \\
                              0           & 0           & \varepsilon
                          \end{bmatrix} = [a_{1}, a_{2}, a_{3}]
    \end{align}
    其中: \\
    \begin{align}
        a_{1} = [1, \varepsilon, 0, 0]^{T} \\
        a_{2} = [1, 0, \varepsilon, 0]^{T} \\
        a_{3} = [1, 0, 0, \varepsilon]^{T}
    \end{align}
    只要: \\
    \begin{align}
        \varepsilon \neq 0
    \end{align} \\
    三列依然线性无关 : \\
    \begin{align}
        rank(A_{\varepsilon}) = 3
    \end{align}
    当 $\varepsilon$ 很小, 三列都非常接近: \\
    \begin{align}
        e_{1} = [1, 0 ,0 ,0]^{T}
    \end{align}
    它们几乎平行 \\

    2: 几何的病态程度: \\

    1.1 长度: \\
    \begin{align}
        ||a_{i}||^{2} = 1 + \varepsilon^{2}
    \end{align}

    1.2 列之间的位置关系(叉积): \\
    \begin{align}
        a_{i}^{T}a_{j} = 1  (i \neq j)
    \end{align}
    不正交 是锐角

    1.3 角度 : \\
    \begin{align}
        \cos\theta  = \frac{<a_{i},a_{j}>}{||a_{i}||||a_{j}||}                   \\
        = \frac{1}{\sqrt{1^{2} + \varepsilon^{2}}\sqrt{1^{2} + \varepsilon^{2}}} \\
        = \frac{1}{1 +  \varepsilon^{2}}                                         \\
    \end{align} \\
    当 $\varepsilon \to 0$ 我们会得到 $\cos\theta \to 1,  \theta \to 0$

    所以我们从几何上就可以发现 当$\varepsilon \to 0$, $a_{i}, a_{2}, a_{3}$ 越接近平行(角度为0 ) \\


    3: Gram matrix 暴露出的结构
    在 传统的正规方程法(Normal Equation), 求解最小二乘问题 (Least Squares Problem) 时 \\
    关键代数式: \\
    \begin{align}
        A^{T}A\hat{c} = A^{T}u \\
        (A^{T}A)\hat{c} = A^{T}u
    \end{align}
    而 $(A^{T}A)$ 就是 Gram matrix

    Gram matrix 可以表现出几个关键的性质: \\
    1: 列与列之间的正交程度 (其他元素与对角线元素越解决 越平行) \\
    2: 因为存在列元素平方的运算, 在计算$\varepsilon$ 这种存在误差的输入时 会指数级拉大误差 \\

    这里我们仔细思考并研究下2性质 : \\
    \begin{align}
        A_{\varepsilon}^{T}A =
        \begin{bmatrix}
            1 & \varepsilon & 0           & 0           \\
            1 & 0           & \varepsilon & 0           \\
            1 & 0           & 0           & \varepsilon \\
        \end{bmatrix}\begin{bmatrix}
                         1           & 1           & 1           \\
                         \varepsilon & 0           & 0           \\
                         0           & \varepsilon & 0           \\
                         0           & 0           & \varepsilon
                     \end{bmatrix}
        = \begin{bmatrix}
              1 + \varepsilon^{2} & 1                   & 1                   \\
              1                   & 1 + \varepsilon^{2} & 1                   \\
              1                   & 1                   & 1 + \varepsilon^{2}
          \end{bmatrix}
    \end{align} \\
    这里我们通过 Gram matrix 会发现 A 列之间的关系是几乎平行的 \\
    我们继续分解: \\
    \begin{align}
        \begin{bmatrix}
            1 + \varepsilon^{2} & 1                   & 1                   \\
            1                   & 1 + \varepsilon^{2} & 1                   \\
            1                   & 1                   & 1 + \varepsilon^{2}
        \end{bmatrix} = \begin{bmatrix}
                            1 & 1 & 1 \\
                            1 & 1 & 1 \\
                            1 & 1 & 1
                        \end{bmatrix} + \begin{bmatrix}
                                            1 & 0 & 0 \\
                                            0 & 1 & 0 \\
                                            0 & 0 & 1
                                        \end{bmatrix}\varepsilon^{2}
    \end{align}
    那么我们能得到 $\begin{bmatrix}
            1 & 1 & 1 \\
            1 & 1 & 1 \\
            1 & 1 & 1
        \end{bmatrix}$ 的 rank = 1, 特征值是$[3, 0, 0] + [\varepsilon^{2},\varepsilon^{2},\varepsilon^{2}]$ \\
    = $[3 + \varepsilon^{2}, \varepsilon^{2}, \varepsilon^{2}]$ \\
    这个特征值组是属于$A_{\varepsilon}^{T}A$ \\
    那么 $A_{\varepsilon}$ 的奇异值呢(对$A_{\varepsilon}^{T}A$ 特征组开平方)? (这里我们还没学习到奇异值, 都可以得到现有公式):\\
    \begin{align}
        \sigma = [\sqrt{3 + \varepsilon^{2}}, \varepsilon, \varepsilon]
    \end{align}
    (奇异值有一个很重要概念: $max(\sigma)$ 代表矩阵在最长方向上把空间放大了多少倍, 对于奇异值而言, 作用的单位并不是一个单个向量,
    而是多个向量(组), 它是可以表示一直几何形状的, 最长方向上把空间放大了多少倍 指的是变化成新的几何形状 方向最长的那个向量. \\
    $min(\sigma)$  代表矩阵在最短方向上把空间缩小了多少倍 \\
    空间就是这个几何形状提供的概念  2d 是面积 3d是体积 ....) \\
    那么我们可以算出: $Condition Number k_{2}$ (误差放大器)\\
    \begin{align}
        K_{2}(A_{\varepsilon}) = \frac{max(\sigma)}{min(\sigma)} = \frac{\sqrt{3 + \varepsilon^{2}}}{\varepsilon}
    \end{align}
    那么 $A_{\varepsilon}^{T}A$ 的奇异值组呢 ?
    这里我们可以先采用现有定理: 对于任意实对称矩阵（并且特征值全部非负），它的奇异值等于它的特征值 \\
    \begin{align}
        K_{2}(A_{\varepsilon}^{T}A) = \frac{max(\sigma)}{min(\sigma)} = \frac{3 + \varepsilon^{2}}{\varepsilon^{2}}
    \end{align}

    所以我么能发现一个很有意思的现象:  $ K_{2}(A_{\varepsilon}^{T}A) ≈ K_{2}(A_{\varepsilon})^{2}$ 是一种平方指数的误差放大 \\
    这正是 normal equation 数值上危险的原因
    这里面还是有一个有意思的问题: 为什么$Condition Number k_{2}$ 是误差放大器 是什么性质导致的(我们还没学习到奇异值).







\end{flushleft}


\end{document}